% Основные источники:
% https://sochisirius.ru/uploads/2021/11/math_1021_11a_12_gaussovy_celye_chisla_teoriya.pdf
% https://sochisirius.ru/uploads/2021/11/math_1021_11a_13_gaussovy_celye_chisla_zadachi.pdf
% https://sochisirius.ru/uploads/2021/11/math_1021_11b_19-celye_gaussovye_chisla.pdf
% https://sochisirius.ru/uploads/2021/11/math_1021_11v_14_gaussovy_chisla.pdf

\begin{center}
    {\LARGE\bfseries Целые гауссовы числа}
\end{center}


\begin{defbox}{гауссовы целые числа и норма}
\emph{Целыми гауссовыми числами} называются комплексные числа вида
\[
    a+bi,\qquad a,b\in\mathbb Z.
\]
Множество всех таких чисел обозначается через
\[
    \mathbb Z[i]=\{a+bi\mid a,b\in\mathbb Z\}.
\]

Для $\alpha=a+bi$ положим
\[
    \overline{\alpha}=a-bi,
    \qquad
    N(\alpha)=\alpha\overline{\alpha}=a^2+b^2.
\]
Число $N(\alpha)$ называется \emph{нормой} числа $\alpha$.
\end{defbox}

\begin{center}
    {\large\bfseries I. Первые свойства}
\end{center}

\problem{1}{Докажите следующие свойства.}

\begin{enumerate}[label=\textbf{\alph*)}, leftmargin=8mm]
    \item Если $\alpha,\beta\in\mathbb Z[i]$, то
    \[
        \alpha+\beta,\quad \alpha-\beta,\quad
        \alpha\beta,\quad \overline{\alpha}
        \in\mathbb Z[i].
    \]

    \item Для любого $\alpha\in\mathbb Z[i]$ число $N(\alpha)$ является
    целым неотрицательным, причём
    \[
        N(\alpha)=0\quad\Longleftrightarrow\quad \alpha=0.
    \]

    \item Для любых $\alpha,\beta\in\mathbb Z[i]$ выполнено
    \[
        N(\alpha\beta)=N(\alpha)N(\beta).
    \]
\end{enumerate}

\smallskip
\textbf{Решение.}
Пусть $\alpha=a+bi$ и $\beta=c+di$, где $a,b,c,d\in\mathbb Z$.

\begin{enumerate}[label=\textbf{\alph*)}, leftmargin=8mm]
    \item Имеем
    \[
    \begin{aligned}
        \alpha+\beta&=(a+c)+(b+d)i,\\
        \alpha-\beta&=(a-c)+(b-d)i,\\
        \alpha\beta&=(ac-bd)+(ad+bc)i,\\
        \overline\alpha&=a-bi.
    \end{aligned}
    \]
    Все действительные и мнимые части здесь целые, поэтому полученные
    числа принадлежат $\mathbb Z[i]$.

    \item $N(\alpha)=a^2+b^2$ является целым неотрицательным числом.
    Равенство $a^2+b^2=0$ возможно ровно тогда, когда $a=b=0$, то есть
    когда $\alpha=0$.

    \item Сопряжение согласовано с умножением:
    $\overline{\alpha\beta}=\overline\alpha\,\overline\beta$. Поэтому
    \[
        N(\alpha\beta)
        =\alpha\beta\overline{\alpha\beta}
        =\alpha\overline\alpha\,\beta\overline\beta
        =N(\alpha)N(\beta).
    \]
\end{enumerate}

\medskip

\begin{defbox}{делимость, обратимые и ассоциированные элементы}
Для $\alpha,\beta\in\mathbb Z[i]$ говорят, что $\alpha$ делит $\beta$,
и пишут $\alpha\mid\beta$, если существует $\gamma\in\mathbb Z[i]$,
для которого
\[
    \beta=\alpha\gamma.
\]

Число $\varepsilon\in\mathbb Z[i]$ называется \emph{обратимым}, если
$\varepsilon^{-1}\in\mathbb Z[i]$.

Числа $\alpha$ и $\beta$ называются \emph{ассоциированными}, если
$\alpha=\varepsilon\beta$ для некоторого обратимого
$\varepsilon\in\mathbb Z[i]$.
\end{defbox}

\problem{2}{Исследуйте обратимые и ассоциированные элементы.}

\begin{enumerate}[label=\textbf{\alph*)}, leftmargin=8mm]
    \item Докажите, что $\varepsilon\in\mathbb Z[i]$ обратимо тогда
    и только тогда, когда
    \[
        N(\varepsilon)=1.
    \]

    \item Найдите все обратимые элементы кольца $\mathbb Z[i]$.

    \item Опишите геометрически все числа, ассоциированные с
    $\alpha=a+bi$.

    \item Верно ли, что два гауссовых числа одинаковой нормы обязательно
    ассоциированы?
\end{enumerate}

\smallskip
\textbf{Решение.}

\begin{enumerate}[label=\textbf{\alph*)}, leftmargin=8mm]
    \item Если $\varepsilon\eta=1$ для некоторого
    $\eta\in\mathbb Z[i]$, то
    \[
        N(\varepsilon)N(\eta)=N(1)=1.
    \]
    Оба множителя являются неотрицательными целыми числами, поэтому
    $N(\varepsilon)=1$. Обратно, при $N(\varepsilon)=1$ выполнено
    $\varepsilon\overline\varepsilon=1$, так что
    $\varepsilon^{-1}=\overline\varepsilon\in\mathbb Z[i]$.

    \item Из $a^2+b^2=1$ получаем
    \[
        \varepsilon\in\{1,-1,i,-i\}.
    \]
    Все четыре числа действительно обратимы.

    \item Все числа, ассоциированные с $\alpha$, имеют вид
    \[
        \alpha,\quad -\alpha,\quad i\alpha,\quad -i\alpha.
    \]
    На комплексной плоскости это точки, полученные из $\alpha$
    поворотами вокруг начала координат на углы
    $0^\circ,90^\circ,180^\circ,270^\circ$.

\begin{center}
\begin{asy}
size(9cm);

pair O=(0,0);
pair a=(2,1);
pair ia=(-1,2);
pair ma=(-2,-1);
pair mia=(1,-2);

real M=2.7;
real R=sqrt(5);

// Координатная сетка
for (int k=-2; k<=2; ++k) {
    draw((k,-M)--(k,M), lightgray+linewidth(0.35));
    draw((-M,k)--(M,k), lightgray+linewidth(0.35));
}

// Координатные оси
draw((-M,0)--(M,0), black+linewidth(0.9), Arrow(TeXHead));
draw((0,-M)--(0,M), black+linewidth(0.9), Arrow(TeXHead));

label("$\operatorname{Re} z$", (M,0), SE);
label("$\operatorname{Im} z$", (0,M), NW);
label("$0$", O, SW);

// Числовые отметки
for (int k=-2; k<=2; ++k) {
    if (k != 0) {
        label(string(k), (k,0), S);
        label(string(k), (0,k), W);
    }
}

// Окружность равных модулей
draw(circle(O,R), gray+dashed+linewidth(0.7));

// Квадрат
draw(
    a--ia--ma--mia--cycle,
    black+linewidth(1.2)
);

// Вершины квадрата
fill(circle(a,0.055), red);
fill(circle(ia,0.055), red);
fill(circle(ma,0.055), red);
fill(circle(mia,0.055), red);

// Подписи
label("$\alpha$", a, NE);
label("$i\alpha$", ia, NW);
label("$-\alpha$", ma, SW);
label("$-i\alpha$", mia, SE);
\end{asy}
\end{center}

    \item Неверно. Например,
    \[
        N(1+2i)=N(2+i)=5,
    \]
    но среди чисел
    \[
        1+2i,\quad -1-2i,\quad -2+i,\quad 2-i
    \]
    число $2+i$ отсутствует. Следовательно, $1+2i$ и $2+i$ не
    ассоциированы.
\end{enumerate}

\medskip

\begin{center}
    {\large\bfseries II. Деление с остатком и алгоритм Евклида}
\end{center}

\problem{3}{
Пусть $z=x+yi$ — произвольное комплексное число. Докажите, что
существует $q=m+ni\in\mathbb Z[i]$, для которого
\[
    |z-q|^2\leqslant\frac12.
\]
}

\textit{Указание.} Выберите $m$ и $n$ ближайшими целыми числами к
$x$ и $y$ соответственно.

\smallskip
\textbf{Решение.}
Выберем целые $m$ и $n$ так, чтобы
\[
    |x-m|\leqslant\frac12,
    \qquad
    |y-n|\leqslant\frac12,
\]
и положим $q=m+ni$. Тогда
\[
    |z-q|^2=(x-m)^2+(y-n)^2
    \leqslant\frac14+\frac14=\frac12.
\]

Геометрически мы выбираем ближайший к $z$ узел квадратной решётки
$\mathbb Z[i]$. Любая точка единичного квадрата находится от его центра
на расстоянии не больше $\sqrt2/2$.

\begin{center}
\begin{asy}
size(8cm);

pair A=(0,0);
pair B=(1,0);
pair C=(1,1);
pair D=(0,1);

pair q=B;
pair z=(0.76,0.22);

draw(A--B--C--D--cycle, black+linewidth(1));
draw(z--q, black+dotted+linewidth(1));

fill(circle(z,0.018), red);
fill(circle(q,0.018), red);

label("$z=x+yi$", z, NW);
label("$q=m+ni$", q, SE);

label(
    "$|z-q|\leqslant\dfrac{\sqrt{2}}{2}$",
    (0.5,-0.20),
    S
);
\end{asy}
\end{center}

\medskip

\problem{4}{
Докажите теорему о делении с остатком: для любых
$\alpha,\beta\in\mathbb Z[i]$, где $\beta\ne0$, существуют
$q,r\in\mathbb Z[i]$ такие, что
\[
    \alpha=\beta q+r,
    \qquad
    N(r)<N(\beta).
\]
}

\textit{Указание.} Примените предыдущую задачу к комплексному числу
$\alpha/\beta$.

\smallskip
\textbf{Решение.}
По задаче 3 для числа $z=\alpha/\beta$ можно выбрать
$q\in\mathbb Z[i]$ так, что
\[
    \left|\frac{\alpha}{\beta}-q\right|^2\leqslant\frac12.
\]
Положим $r=\alpha-\beta q$. Тогда $r\in\mathbb Z[i]$ и
$\alpha=\beta q+r$. Кроме того,
\[
\begin{aligned}
    N(r)=|r|^2
    &=|\beta|^2
      \left|\frac{\alpha}{\beta}-q\right|^2\\
    &\leqslant\frac12|\beta|^2
     =\frac12N(\beta)<N(\beta).
\end{aligned}
\]
Это и есть требуемое деление с остатком.

\medskip

\begin{remarkbox}
В отличие от деления целых чисел, неполное частное и остаток
в $\mathbb Z[i]$ могут определяться неоднозначно.
\end{remarkbox}

\begin{defbox}{наибольший общий делитель}
Гауссово число $d$ называется \emph{наибольшим общим делителем}
чисел $\alpha$ и $\beta$, если
\[
    d\mid\alpha,\qquad d\mid\beta,
\]
и каждый общий делитель $\alpha$ и $\beta$ делит $d$.

Наибольший общий делитель определён с точностью до умножения на
обратимый элемент.
\end{defbox}

\problem{5}{Исследуйте алгоритм Евклида в $\mathbb Z[i]$.}

\begin{enumerate}[label=\textbf{\alph*)}, leftmargin=8mm]
    \item Докажите, что последовательное деление с остатком позволяет
    найти наибольший общий делитель двух гауссовых чисел.

    \item Найдите
    \[
        \text{НОД}(4+7i,\;9+2i).
    \]

    \item Представьте найденный наибольший общий делитель в виде
    \[
        (4+7i)u+(9+2i)v,
        \qquad u,v\in\mathbb Z[i].
    \]

    \item Найдите
    \[
        \text{НОД}(3+8i,\;7-5i).
    \]
\end{enumerate}

\smallskip
\textbf{Решение.}

\begin{enumerate}[label=\textbf{\alph*)}, leftmargin=8mm]
    \item Последовательно делим с остатком:
    \[
    \begin{aligned}
        \alpha&=\beta q_1+r_1,\\
        \beta&=r_1q_2+r_2,\\
        r_1&=r_2q_3+r_3,\\
        &\ \vdots
    \end{aligned}
    \qquad
    N(\beta)>N(r_1)>N(r_2)>\dotsb.
    \]
    Нормы являются неотрицательными целыми числами, поэтому процесс
    заканчивается. Пусть $r_s$~--- последний ненулевой остаток. Из
    последнего равенства следует, что $r_s$ делит предыдущий остаток;
    поднимаясь по равенствам вверх, получаем
    $r_s\mid\alpha$ и $r_s\mid\beta$.

    Наоборот, любой общий делитель $\alpha$ и $\beta$ делит $r_1$,
    затем $r_2$ и так далее, а значит, делит $r_s$. Следовательно,
    $r_s$ является наибольшим общим делителем.

    \item Делим $9+2i$ на $4+7i$:
    \[
        9+2i=(4+7i)(1-i)+(-2-i).
    \]
    Далее деление выполняется без остатка:
    \[
        4+7i=(-2-i)(-3-2i).
    \]
    Поэтому
    \[
        \text{НОД}(4+7i,9+2i)=-2-i
    \]
    с точностью до умножения на $\pm1,\pm i$.

    \item Из первого равенства сразу получаем
    \[
        -2-i=(4+7i)(-1+i)+(9+2i)\cdot1.
    \]
    Таким образом, можно взять
    \[
        u=-1+i,\qquad v=1.
    \]

    \item Имеем
    \[
    \begin{aligned}
        7-5i&=(3+8i)(-i)+(-1-2i),\\
        3+8i&=(-1-2i)(-4)-1.
    \end{aligned}
    \]
    Последний ненулевой остаток равен $-1$, то есть обратим. Поэтому
    \[
        \text{НОД}(3+8i,7-5i)=1
    \]
    с точностью до ассоциированности.
\end{enumerate}

\newpage

\begin{defbox}{простое гауссово число}
Ненулевое необратимое число $\pi\in\mathbb Z[i]$ называется
\emph{простым гауссовым числом}, если из равенства
\[
    \pi=\alpha\beta
\]
следует, что хотя бы одно из чисел $\alpha,\beta$ обратимо.
\end{defbox}

\problem{6}{Докажите аналоги стандартных арифметических фактов.}

\begin{enumerate}[label=\textbf{\alph*)}, leftmargin=8mm]
    \item Наибольший общий делитель любых $\alpha,\beta\in\mathbb Z[i]$
    можно представить в виде
    \[
        d=\alpha u+\beta v,
        \qquad u,v\in\mathbb Z[i].
    \]

    \item Если $\text{НОД}(\alpha,\beta)=1$ и
    \[
        \alpha\mid\beta\gamma,
    \]
    то $\alpha\mid\gamma$.

    \item Если $\pi$ — простое гауссово число и
    \[
        \pi\mid\alpha\beta,
    \]
    то $\pi\mid\alpha$ или $\pi\mid\beta$.

    \item Докажите существование разложения любого ненулевого
    необратимого гауссова числа в произведение простых.

    \item Докажите единственность такого разложения с точностью до
    порядка сомножителей и замены сомножителей ассоциированными.
\end{enumerate}

\smallskip
\textbf{Решение.}

\begin{enumerate}[label=\textbf{\alph*)}, leftmargin=8mm]
    \item В алгоритме Евклида каждый остаток является линейной
    комбинацией двух предыдущих. Поэтому, последовательно выражая
    остатки через предыдущие, последний ненулевой остаток $d$ получаем
    в виде
    \[
        d=\alpha u+\beta v,
        \qquad u,v\in\mathbb Z[i].
    \]

    \item Из пункта а) существуют $u,v\in\mathbb Z[i]$ такие, что
    \[
        \alpha u+\beta v=1.
    \]
    Умножим равенство на $\gamma$:
    \[
        \gamma=\alpha(\gamma u)+(\beta\gamma)v.
    \]
    Оба слагаемых справа делятся на $\alpha$: первое очевидно, а второе
    по условию. Значит, $\alpha\mid\gamma$.

    \item Предположим, что $\pi\nmid\alpha$. Любой общий делитель
    $\pi$ и $\alpha$ делит простое число $\pi$, поэтому он либо обратим,
    либо ассоциирован с $\pi$. Второй случай невозможен, иначе
    $\pi\mid\alpha$. Следовательно,
    \[
        \text{НОД}(\pi,\alpha)=1.
    \]
    Из $\pi\mid\alpha\beta$ и пункта б) получаем $\pi\mid\beta$.

    \item Проведём индукцию по $N(\alpha)$. Если $\alpha$ простое,
    разложение уже найдено. Если нет, то
    \[
        \alpha=\beta\gamma,
    \]
    где $\beta$ и $\gamma$ необратимы. Поэтому
    \[
        1<N(\beta)<N(\alpha),
        \qquad
        1<N(\gamma)<N(\alpha).
    \]
    По предположению индукции оба множителя раскладываются на простые,
    а значит, раскладывается и $\alpha$.

    \item Пусть получены два разложения
    \[
        \pi_1\pi_2\dotsm\pi_r
        =\rho_1\rho_2\dotsm\rho_s.
    \]
    Число $\pi_1$ делит правую часть, поэтому по пункту в) делит один
    из множителей $\rho_j$. Так как $\rho_j$ простое, числа $\pi_1$ и
    $\rho_j$ ассоциированы. Сократим их и повторим рассуждение.
    В результате все простые множители двух разложений попарно
    сопоставятся с точностью до ассоциированности; в частности,
    $r=s$. Единственность доказана.
\end{enumerate}

\begin{factbox}{основная теорема арифметики в $\mathbb Z[i]$}
Любое ненулевое гауссово число раскладывается в произведение простых
гауссовых чисел. Такое разложение единственно с точностью до порядка
сомножителей и умножения сомножителей на обратимые элементы.

Этим фактом можно пользоваться во всех последующих задачах.
\end{factbox}

\begin{center}
    {\large\bfseries III. Какие гауссовы числа являются простыми?}
\end{center}

\problem{7}{
Докажите, что если $N(\alpha)$ является обычным простым числом,
то $\alpha$ является простым гауссовым числом.
}

\smallskip
\textbf{Решение.}
Пусть $N(\alpha)=p$, где $p$~--- обычное простое число, и
$\alpha=\beta\gamma$. Тогда
\[
    p=N(\alpha)=N(\beta)N(\gamma).
\]
Так как $N(\beta)$ и $N(\gamma)$ являются неотрицательными целыми
числами, одно из них равно $1$. Соответствующий множитель обратим по
задаче 2. Следовательно, $\alpha$ является простым гауссовым числом.

\medskip

\problem{8}{
Пусть $p$ — обычное простое число вида $4k+3$. Докажите, что
$p$ является простым числом и в кольце $\mathbb Z[i]$.
}

\textit{Указание.} Если $p=\alpha\beta$, рассмотрите возможные значения
$N(\alpha)$ и $N(\beta)$ и вспомните остатки квадратов по модулю $4$.

\smallskip
\textbf{Решение.}
Предположим, что $p=\alpha\beta$, причём оба множителя необратимы.
После перехода к нормам получаем
\[
    p^2=N(p)=N(\alpha)N(\beta).
\]
Нормы обоих множителей больше $1$, поэтому из простоты $p$ следует
\[
    N(\alpha)=N(\beta)=p.
\]
Если $\alpha=a+bi$, то отсюда
\[
    p=a^2+b^2.
\]
Но квадрат целого числа сравним по модулю $4$ только с $0$ или $1$,
поэтому сумма двух квадратов не может быть сравнима с $3$ по модулю
$4$. Противоречие. Значит, в любом разложении $p=\alpha\beta$
один из множителей обратим, то есть $p$ прост в $\mathbb Z[i]$.